% Figure 1 (한글판): 연구 구조. 단순 블록 배치. 색·점선 없음, 화살표 최소.
\begin{figure*}[t]
\centering
\resizebox{\textwidth}{!}{%
\begin{tikzpicture}[
  font=\sffamily,
  cell/.style={rectangle, draw=black, minimum height=1.15cm, minimum width=4.0cm,
               align=center, inner sep=4pt},
  band/.style={rectangle, draw=black, thick, minimum height=1.15cm, minimum width=16.6cm,
               align=center, inner sep=4pt},
  arr/.style={-{Stealth[length=6pt]}, thick},
]

% 1행: 우리가 바꾸는 네 개의 축
\node[cell] (a1) at (0,0)     {\textbf{모델}\\{\small 상용 8종, 오픈 4종}};
\node[cell] (a2) at (4.2,0)   {\textbf{하드웨어}\\{\small RTX PRO 6000, H100, A6000}};
\node[cell] (a3) at (8.4,0)   {\textbf{벤치마크}\\{\small 4종, 424개 과제}};
\node[cell] (a4) at (12.6,0)  {\textbf{예산}\\{\small 샘플 수, 라운드 수, 턴 수}};

% 2행: 고정된 하니스
\node[band] (harness) at (6.3,-2.15)
  {\textbf{하나의 고정된 평가 하니스}\\[2pt]
   {\small 프롬프트 구성 \ \textbar\ \ 생성 \ \textbar\ \ 코드 추출 \ \textbar\ \
    \textbf{격리된 서브프로세스에서 채점} \ \textbar\ \ 에러 분류 \ \textbar\ \ 수정과 재시도}};

% 3행: 관측량
\node[band] (out) at (6.3,-4.15)
  {\textbf{측정 대상}\\[2pt]
   {\small 정확성, Fast-0과 Pass@1 \ \textbar\ \ 속도, Fast-1과 Fast-2 \ \textbar\ \
    에러 유형 \ \textbar\ \ 벽시계 시간과 완주율}};

% 4행: 연구 질문
\node[band] (rq) at (6.3,-6.15)
  {\textbf{연구 질문}\\[2pt]
   {\small RQ1 얼마나 정확한가 \ \textbar\ \ RQ2 설정이 얼마나 좌우하는가 \ \textbar\ \
    RQ3 실패는 어떤 모습인가 \ \textbar\ \ RQ4 무엇이 정답률을 올리는가 \ \textbar\ \
    RQ5 무엇이 이식되는가}};

% 화살표 세 개만
\draw[arr] (6.3,-0.62) -- (harness.north);
\draw[arr] (harness.south) -- (out.north);
\draw[arr] (out.south) -- (rq.north);

\end{tikzpicture}%
}
\caption{본 연구의 구조. 하나의 평가 하니스를 고정하고 그 위에서 네 개의 축을 바꾼다.
하니스는 모든 커널을 격리된 서브프로세스에서 채점한다. 이는 선택이 아니라 필수인데,
잘못된 커널 하나가 CUDA 컨텍스트를 오염시키면 같은 프로세스의 이후 평가가 모두 실패하기 때문이다.
우리는 새로운 생성 시스템을 제안하지 않는다. 하니스를 고정한 채 한 번에 하나의 축만 바꾸어
모델, 하드웨어, 벤치마크, 계산 예산의 효과를 분리한다.}
\label{fig:overview}
\end{figure*}
